\section{What the Composer Is Actually Checking}
\label{sec:ch16-what-the-composer-is-checking}

\index{composition!what it is asking, and what it is not|(}%
Three edits composed cleanly in
section~\ref{sec:ch16-three-changes-and-one-status-code}, which invites the
reading that the composer is lax. It is not. It is answering a question, the
question is a good one, and it is not this one. The way to see that is to give
it fourteen edits and watch which two it stops.

\index{composition!fourteen edits, one at a time}%
Each row below is one change to one document, with the other three left
exactly as the services export them, composed as a four-subgraph graph. The
third column is where that change sits on the list
section~\ref{sec:ch16-the-change-this-book-already-shipped} took out of
\code{graphql-js}.

\begin{center}
\footnotesize
\begin{tabular}{l l l}
\toprule
Change to one document & Composition & On the list as \tabularnewline
\midrule
remove \code{Query.sessionById} & exit 0 & \code{FIELD\_REMOVED}
  \tabularnewline
remove \code{Session.abstract} & exit 0 & \code{FIELD\_REMOVED}
  \tabularnewline
rename \code{Session.title} to \code{name} & \textbf{exit 1} &
  \code{FIELD\_REMOVED} \tabularnewline
\code{Session.title} loses its \code{!} & exit 0 & \code{FIELD\_CHANGED\_KIND}
  \tabularnewline
\code{Session.durationMinutes} loses its \code{!} & exit 0 &
  \code{FIELD\_CHANGED\_KIND} \tabularnewline
\code{sessions} gains \code{track: String!} & exit 0 &
  \code{REQUIRED\_ARG\_ADDED} \tabularnewline
\code{sessionById(id: Int!)} becomes \code{ID!} & exit 0 &
  \code{ARG\_CHANGED\_KIND} \tabularnewline
remove \code{Session.startsAt} & exit 0 & \code{FIELD\_REMOVED}
  \tabularnewline
widen the \code{@key} to \code{"id title"} & \textbf{exit 1} & not on it
  \tabularnewline
remove \code{Speaker.bio} & exit 0 & \code{FIELD\_REMOVED} \tabularnewline
\code{Speaker.name} loses its \code{!} & exit 0 & \code{FIELD\_CHANGED\_KIND}
  \tabularnewline
\code{Session.averageScore} becomes \code{Int} & exit 0 &
  \code{FIELD\_CHANGED\_KIND} \tabularnewline
remove \code{Session.ratingCount} & exit 0 & \code{FIELD\_REMOVED}
  \tabularnewline
remove \code{Query.searchSessions} & exit 0 & \code{FIELD\_REMOVED}
  \tabularnewline
\bottomrule
\end{tabular}
\end{center}

\index{composition!twelve of fourteen say nothing}%
Twelve of the fourteen compose, and each of the twelve printed the success
line and nothing else. Every one of those twelve is on the breaking list.

\index{composition!the two it refuses are not the two you would guess}%
The two it refuses are on the other side of the ledger. One of them, the
rename, is on the list as well. The other, widening a key, is not on it at
all: a \code{@key} is a federation directive rather than something a client
can select, and no client will ever know the difference. So the composer
stopped one change no client can see and passed twelve that break one.

\index{composition!the row worth arguing about}%
One row is worth arguing about and
section~\ref{sec:ch16-the-gate-and-what-it-needs} has the vendors' own answer
to the argument. \code{sessionById(id: Int!)} becoming \code{id: ID!} is
\code{ARG\_CHANGED\_KIND} because the two named types differ, and a reader who
knows that \code{ID} accepts an integer will want to say it is harmless. A
schema check will still call it breaking, and Apollo's documentation says
plainly why: validation does not have the information to establish that it is
not.

\subsection{Why Those Two, and Not the Reasons You Would Predict}

\index{composition!the rename fails at the far end}%
Both refusals are worth reading, because neither fails where you would look.
Rename \code{Session.title} in the Sessions document and the message is not
about the rename:

\begin{minted}[fontsize=\footnotesize,breakanywhere]{text}
The Object "Session" is invalid because the following field definition is declared "@external" on all instances of
that field:
 "title" in subgraph "ratings"
At least one instance of a field definition must always be resolvable (and therefore not declared "@external").
\end{minted}

\index{external@\code{"@external}!the last declaration standing}%
Chapter~\ref{ch:second-third-subgraph} gave the Ratings document a copy of
\code{title} marked \code{@external}, so that \code{feedbackUrl} could
\code{@requires} it.
Rename the Sessions copy and that borrowed declaration becomes the only
declaration of a field by that name anywhere, which the composer refuses
categorically. It is not complaining about the \code{@requires} selection set
and it is not complaining about Sessions. It is complaining that a field
exists which nobody can resolve, in the subgraph that was only ever borrowing
it.

\index{key@\code{"@key}!widening it cuts somebody else's route}%
Widening the key fails somewhere else again. Add \code{title} to the Sessions
\code{@key} and Ratings still declares \code{@key(fields: "id")} on its own
\code{Session}, so it can no longer satisfy the stricter key and the route
Search had into Sessions through Ratings is cut. Four errors come back, one
per orphaned field, in exactly the shape
chapter~\ref{ch:satisfiability} spent a chapter reading: \code{speaker},
\code{abstract}, \code{startsAt} and \code{durationMinutes}, all named at a
path through \code{searchSessions}, with \code{ratings} as the entity ancestor.

\index{composition!both refusals are about routing}%
So the two refusals are one refusal, arrived at twice: the composer could not
plan a route, which is the only question it was ever asked.

\subsection{Every Question It Can Ask Is a Question About the Files}

\index{composition!reads four documents and nothing else}%
Chapter~\ref{ch:composition} described composition as two passes, a merge and
a walk, and chapter~\ref{ch:satisfiability} spent its length on the second.
What both passes have in common is their input. Four documents go in. Nothing
else does. \code{wgc router compose} does not call a service, which
chapter~\ref{ch:composition} treated as a feature and still is, and it does
not read a client, which nobody has ever offered as a feature because nobody
has ever offered it at all.

\index{composition!a client operation is in no file}%
So every question it can answer is a question the four files answer. Can these
types merge. Does every field have somewhere it can be resolved from. Is there
a path from each root field to each leaf. \emph{Was this field being used} is
not that kind of question. The operation that uses it is in a mobile
application somebody shipped a while ago, or in a dashboard in another
repository, or in a script written once and scheduled. It is not in the four
documents, it is not in the repository the four documents live in, and no
amount of reading them harder will produce it.

\index{figure!where a change fails}%
Figure~\ref{fig:ch16-two-places-a-change-can-fail} is the two outcomes on one
timeline.

\begin{figure}[htbp]
  \centering
  % The same shape of edit to the sessions subgraph, twice: a renamed field the
% composer refuses at compose time, and a removed field it lets straight
% through. The point is where each failure lands on the timeline, not what
% either edit looks like, which is why this is a timeline and not a box
% diagram.
%
% Same idiom as figures/tikz/ch01-one-field.tex, which fixes it: two lanes,
% x=1mm/y=1mm, ticks above a single-weight axis, event nodes anchored south,
% a span bracket under the stage a label needs to explain. The outcome node
% at the right of each lane is figures/tikz/ch15-the-order-decides.tex's
% outcome style, including that file's choice to place both lanes' outcome
% nodes at the same x position even though lane A's axis stops well short of
% it: the blank run-in says as much as the text does.
%
% The cross mark on both lanes is figures/tikz/ch07-the-walk.tex's mark for a
% step that goes wrong at that point. It is used a little differently here
% than in ch07 or in ch15's hollow square: neither lane draws a dashed
% lead-in. Dashing already means, in both of those figures, a stage nobody
% actually walked. Every stage in this figure did happen; the composer really
% ran, the request really went out. So the axis stays solid right up to the
% cross in both lanes, and only the cross mark itself, not the line beneath
% it, carries the meaning. Lane A's axis simply ends there (house style: draw
% only the stages a run actually reaches), while lane B's runs on past its
% cross to the same right margin as lane A's outcome column, because the
% request that fails is still the last event on that lane, not an aborted one.
%
% No quotation marks anywhere in this file, including this comment block, per
% the note at the head of ch07-the-walk.tex and repeated in
% ch15-the-order-decides.tex.
%
% Bare tikzpicture (house style): the figure environment, caption and label
% live at the call site in the section file.
\begin{tikzpicture}[
  x=1mm, y=1mm,
  every node/.append style={font=\sffamily\scriptsize},
  axis/.style={draw, line width=0.4pt,
    -{Straight Barb[length=1.4mm, width=1.6mm]}},
  tick/.style={draw, line width=0.4pt},
  event/.style={anchor=south, align=center, text width=30mm, inner sep=1pt},
  outcome/.style={anchor=west, align=left, text width=32mm, inner sep=1pt},
  span/.style={draw, line width=0.4pt},
  spanlabel/.style={anchor=north, align=center, inner sep=2pt},
  lane/.style={anchor=west, font=\sffamily\scriptsize\itshape},
]

% ------------------------------------------------------------------- lane A
% Both lane titles sit at 17 rather than at ch01's 13, because the first
% event in each lane runs to four lines where ch01's ran to two, and at 13
% the title and the label collide.
\node[lane] at (0,17) {A. a change the composer refuses};

% Solid the whole way: the composer does run compose, and fails there. The
% axis stops short of the cross so that its arrowhead points into the cross
% rather than being drawn underneath it, and nothing is drawn past it.
\draw[axis] (0,0) -- (41,0);
\draw[tick] (14,-1.5) -- (14,1.5);
\draw[tick] (42.8,-1.5) -- (45.2,1.5);
\draw[tick] (42.8,1.5) -- (45.2,-1.5);

\node[event, text width=34mm] at (14,3)
  {edit the sessions document\\(rename Session.title to Session.name)};
\node[event, text width=26mm] at (44,3) {wgc router compose\\exit 1};

\draw[span] (7,-4) -- (7,-7) -- (51,-7) -- (51,-4);
\node[spanlabel, text width=40mm] at (29,-8)
  {the composer can answer this: the rename and the read it breaks are both
   in the four documents it composes};

\node[outcome] at (118,0) {refused. Nothing\\deployed.};

% ------------------------------------------------------------------- lane B
\begin{scope}[shift={(0,-50)}]
  \node[lane] at (0,17) {B. the same shape of change, waved through};

  \draw[axis] (0,0) -- (99,0);
  \foreach \x in {14,44,74}{\draw[tick] (\x,-1.5) -- (\x,1.5);}
  \draw[tick] (100.8,-1.5) -- (103.2,1.5);
  \draw[tick] (100.8,1.5) -- (103.2,-1.5);

  % No chapter number in the printed text: the book writes those as a
  % cross-reference, and a figure cannot. It says long deprecated instead,
  % which is what the reader needs here anyway.
  \node[event, text width=34mm] at (14,3)
    {edit the sessions document\\(remove Query.sessionById, long
     deprecated)};
  \node[event, text width=30mm] at (44,3)
    {wgc router compose\\exit 0, nothing printed};
  \node[event, text width=20mm] at (74,3) {the router\\loads it};
  \node[event, text width=26mm] at (102,3)
    {a client sends\\the old query};

  \draw[span] (95,-4) -- (95,-7) -- (109,-7) -- (109,-4);
  \node[spanlabel, text width=34mm] at (102,-8)
    {the operation that fails is in nobody's repository, so no check
     before this point had it to run against};

  \node[outcome] at (118,0) {HTTP 200, and\\an errors array.};
\end{scope}

\end{tikzpicture}

  \caption{The same shape of edit, to the same document, ending in two places.
  Above, the composer refuses and nothing is deployed, because both halves of
  the question it asked were in the documents it was handed. Below, the
  composition passes, the router loads the config, and the failure waits for a
  request the graph had no copy of. The distance between the two crosses is a
  deployment, and everything that makes a breaking change expensive is in that
  gap.}
  \label{fig:ch16-two-places-a-change-can-fail}
\end{figure}

\index{composition!a pass is a pass on one question}%
None of this is an argument against composing. A refusal is real every time,
and chapter~\ref{ch:satisfiability} already said what a pass is worth on its
own terms. What this section adds is narrower and worse: a composer that
answered its own question perfectly would still pass all twelve of the changes
above, because the thing they break was never in the room. Which is the
question the next section takes to the people who sell an answer to it.
\index{composition!what it is asking, and what it is not|)}%
